\chapter[\emj{🔎}~Búsqueda de Patrones: KMP y Rabin-Karp]{\emj{🔎}~Búsqueda de Patrones en Texto: KMP y Rabin-Karp}

\begin{dsaquees}

Buscar una palabra (patrón) dentro de un texto largo 📃, como el
Ctrl+F de tu navegador. La forma ingenua compara letra por letra desde
cada posición: lenta ($O(n \cdot m)$) porque repite comparaciones.

Dos formas inteligentes de hacerlo en $O(n + m)$:
\begin{itemize}
\item KMP: precalcula cuánto puede "saltar" el patrón cuando falla, sin
      retroceder en el texto.
\item Rabin-Karp: convierte cada ventana del texto en un número (hash) y
      compara números en vez de letras.
\end{itemize}

\end{dsaquees}

\begin{dsanino}

Quieres encontrar la palabra "ABABC" dentro de un texto larguísimo 🔤.

La forma tonta: la comparas empezando en cada letra, y si falla, vuelves
a empezar todo desde la siguiente letra. Repites trabajo que ya hiciste.

KMP es más listo: cuando ya comparaste "ABAB" y falla la siguiente letra,
KMP sabe que "AB" del inicio del patrón ya coincide, así que NO vuelve
hasta el principio; salta directo y sigue. Es como cuando ya memorizaste
parte de una canción y no empiezas desde cero cada vez que te equivocas.

Rabin-Karp usa otro truco: le pone un "número de identidad" (huella) a
cada pedazo del texto del tamaño de la palabra, y compara huellas. Si dos
huellas son iguales, probablemente encontró la palabra.

\end{dsanino}

\begin{dsaotraforma}

\textbf{KMP (Knuth-Morris-Pratt)} evita retroceder en el texto usando un
arreglo de "prefijos-sufijos" (\verb|lps|, longest proper prefix which is
also suffix). \verb|lps[i]| dice, para el prefijo del patrón hasta $i$,
cuántos caracteres del inicio coinciden con el final. Cuando una
comparación falla, el patrón salta a \verb|lps[j-1]| en vez de reiniciar.
Total: $O(n + m)$.

\textbf{Rabin-Karp} calcula un hash rodante (rolling hash) de cada
ventana del texto de longitud $m$ y lo compara con el hash del patrón.
El rolling hash actualiza en $O(1)$ al deslizar la ventana (quita la
letra que sale, agrega la que entra). Promedio $O(n + m)$; peor caso
$O(n \cdot m)$ si hay muchas colisiones de hash. Brilla para buscar
MÚLTIPLES patrones a la vez.

\end{dsaotraforma}

\begin{dsadiagrama}
\begin{verbatim}
KMP: patrón = "ABABCABAB"

Construir lps (prefijo propio = sufijo más largo):
  índice:  A  B  A  B  C  A  B  A  B
  lps:     0  0  1  2  0  1  2  3  4

Significado: si fallas en la posición 8 (ya casaste "ABABCABA"),
saltas a lps[7]=3, o sea reusas "ABA" ya emparejado. No retrocedes
en el texto.

RABIN-KARP: rolling hash
  texto: "ABCDE",  patrón "CDE" (largo 3)
  hash("ABC") -> desliza -> hash("BCD") -> hash("CDE") == hash patrón!
  (al deslizar: quita 'A', agrega 'D', todo en O(1))
\end{verbatim}
\end{dsadiagrama}

\begin{lstlisting}
KMP EN 2 FASES
1. Precalcula lps[] del patrón: O(m).
2. Recorre el texto con dos índices (i en texto, j en patrón). Si
   coinciden, avanza ambos. Si falla y j>0, j = lps[j-1] (salta). Si
   j llega a m, encontraste una ocurrencia.

📝 PSEUDOCÓDIGO (construir lps)
──────────────────────────────────────────────────────────────

función buildLPS(patrón):
    lps[0] = 0; len = 0; i = 1
    MIENTRAS i < m:
        SI patrón[i] == patrón[len]:
            len++; lps[i] = len; i++
        SINO SI len > 0:
            len = lps[len-1]          // retrocede en el patrón
        SINO:
            lps[i] = 0; i++
    return lps

💻 CÓDIGO C++ (KMP)
──────────────────────────────────────────────────────────────

#include <string>
#include <vector>
using namespace std;

vector<int> buildLPS(const string& p) {
    int m = p.size(); vector<int> lps(m, 0);
    for (int i = 1, len = 0; i < m; ) {
        if (p[i] == p[len]) lps[i++] = ++len;
        else if (len) len = lps[len-1];
        else lps[i++] = 0;
    }
    return lps;
}

int kmpSearch(const string& text, const string& pat) {
    vector<int> lps = buildLPS(pat);
    int n = text.size(), m = pat.size();
    for (int i = 0, j = 0; i < n; ) {
        if (text[i] == pat[j]) { i++; j++;
            if (j == m) return i - m;      // ocurrencia encontrada
        } else if (j) j = lps[j-1];        // salta usando lps
        else i++;
    }
    return -1;
}

💻 CÓDIGO C++ (Rabin-Karp, rolling hash)
──────────────────────────────────────────────────────────────

int rabinKarp(const string& text, const string& pat) {
    int n = text.size(), m = pat.size();
    if (m > n) return -1;
    const long long B = 256, MOD = 1e9 + 7;
    long long hp = 0, ht = 0, pow = 1;
    for (int i = 0; i < m; i++) {
        hp = (hp * B + pat[i]) % MOD;
        ht = (ht * B + text[i]) % MOD;
        if (i) pow = pow * B % MOD;
    }
    for (int i = 0; i + m <= n; i++) {
        if (hp == ht && text.compare(i, m, pat) == 0) return i; // verifica
        if (i + m < n)                                  // desliza ventana
            ht = ((ht - text[i]*pow % MOD + MOD) * B + text[i+m]) % MOD;
    }
    return -1;
}

⏱️ COMPLEJIDAD (Big-O)
──────────────────────────────────────────────────────────────

  Algoritmo     │ Tiempo          │ Nota
  ──────────────┼─────────────────┼──────────────────────
  Naive         │ O(n·m)          │ compara desde cada pos
  KMP           │ O(n + m)        │ nunca retrocede en texto
  Rabin-Karp    │ O(n + m) prom.  │ O(n·m) peor (colisiones)

* Rabin-Karp brilla buscando múltiples patrones a la vez.
\end{lstlisting}

\begin{dsaentrevistas}

✅ El corazón de KMP es el lps: si sabes construir el arreglo de
   prefijos-sufijos, KMP sale solo. Practica el lps hasta que sea
   automático.

💡 "Implement strStr()" (LC 28) y "Repeated String Match" son KMP
   directos. "Longest Happy Prefix" (LC 1392) es literalmente
   \verb|lps[m-1]|.

⚠️ Rabin-Karp necesita verificación: dos hashes iguales NO garantizan
   igualdad (colisión). Siempre compara los caracteres reales al hacer
   match de hash.

🔑 Rolling hash reutilizable: la técnica de hash rodante de Rabin-Karp
   sirve para detectar substrings duplicados, comparar rangos de texto en
   O(1) y problemas de "longest duplicate substring".

\end{dsaentrevistas}

\begin{dsaresumen}

• Buscar un patrón de largo m en un texto de largo n.
• Naive O(n·m); KMP y Rabin-Karp O(n + m).
• KMP: arreglo lps (prefijo propio = sufijo) para saltar sin retroceder.
• Rabin-Karp: rolling hash de cada ventana; compara números.
• Rabin-Karp exige verificar el match real (colisiones de hash).
• Rabin-Karp es ideal para múltiples patrones a la vez.
════════════════════════════════════════════════════════════════

\end{dsaresumen}
